\section{Estructura de la memoria}\label{cap:estructura_memoria}

Se debe actualizar al terminar el TFG!!!!!!!!!!!!!!!!!!!!!


La memoria se ha organizado en distintos capítulos en los cuales se puede ver la estructura y el desarrollo de dicho proyecto, de forma que se pueda comprender el proyecto de forma clara y completa.

En primer lugar, en este capítulo de introducción se presenta el contexto del proyecto, la motivación que lo impulsa, el estado actual de las distintas soluciones existentes, así como los objetivos que se llevan a cabo.

Después, en el capítulo de planificación se describe la metodología seguida durante el desarrollo de la aplicación, la organización del proyecto, además de las distintas fases de desarrollo y gestión de tareas.

A continuación, en el capítulo de análisis se identifican los actores del sistema, se especifican los requisitos funcionales, de información y no funcionales, y se detallan los casos de uso y el modelo de dominio sobre los que se sustenta el desarrollo de la aplicación.

Seguido por el capítulo de las iteraciones, en el cual se muestran las distintas etapas en las que se ha dividido y repartido el desarrollo de la aplicación, mostrando la evolución que ha ido teniendo el sistema a lo largo del desarrollo del proyecto.

Además, cuenta con un capítulo de estado, en el que se muestra el estado final del sistema desarrollado.

Por último, incluye un capítulo de pruebas en el que se analizan los resultados obtenidos al conseguir construir la versión final y operativa de la aplicación.

Esta memoria concluye con un capítulo de conclusiones, donde se valoran los objetivos conseguidos y se plantean posibles mejoras futuras.